% ============================================================
% Section 9: Conclusion
% ============================================================

\section{Conclusion}
\label{sec:conclusion}

The key-value cache of autoregressive transformers is not merely a memory management structure.
It is a first-class interpretability signal that has been overlooked by the alignment and mechanistic interpretability communities despite being accessible, model-agnostic in form, and practically informative.

We have demonstrated that cognitive mode---metacognitive, analytical, affective, and misalignment states---is readable from \kvcache geometry with accuracy sufficient for deployment as a safety monitoring signal.
Within-model detection of deception, confabulation, sycophancy, and refusal achieves \auroc of 0.78--1.000 after rigorous confound control (0.93--0.995 for the best-controlled comparisons with equal-length system prompts), using four geometric features and logistic regression.
A four-condition sycophancy experiment demonstrates that detection reflects cognitive mode rather than output accuracy: sycophantic and contrarian outputs, which produce the same error rate, are geometrically distinguishable at \auroc 0.757--0.942.

The geometry does more than classify.
It reveals cognitive architecture: confabulation and deception are related but distinct operations---both deviate from honest baselines in the same general direction (expansion on effective rank, contraction on norm-per-token), but confabulation shows larger effects on both dimensions, with a disproportionately strong norm-per-token contraction.
They transfer differently across architectures and engage different dominant features in each model family.
Confabulation---generating beyond the knowledge boundary---is a universal computational challenge that produces transferable geometric signatures, though we note that the cross-model confabulation transfer result (0.93--0.995) requires confirmation with equal-length system prompts, as it may partly reflect a prompt-length artifact.
Deception---suppressing known-correct information---employs architecture-specific circuits that resist cross-model generalization.

A central methodological contribution is the demonstration that rigorous confound control is non-negotiable for \kvcache analysis.
Approximately half of our raw findings across five months of experimentation proved to be token-count artifacts.
Frisch-Waugh-Lovell residualization, same-prompt designs, dual-phase extraction, and encoding-phase baselines are not optional refinements---they are necessary conditions for valid conclusions.
We encourage future work in this space to adopt these controls from the outset rather than discovering confounds post hoc, as we did.

Several open questions remain.
Causal validation through geometry-guided abliteration would determine whether cache geometry participates in cognitive mode or merely reflects it.
Applicability to closed-model APIs requires either provider cooperation or the development of proxy methods that infer geometric properties from observable outputs.
Real-time intervention---modifying cache geometry during generation to steer cognitive mode---is a natural extension that could complement representation engineering with a cache-level control surface.
Scale testing beyond 70B, and characterization of architectures beyond the Qwen/Llama/Mistral family, would establish the generality of these findings across the full landscape of deployed models.

We release our methodology, statistical framework, and aggregate results in full to enable independent replication and critique.
Model-specific direction vectors are staged for responsible disclosure, following the ethical framework detailed in Section~\ref{sec:ethics}.
The \kvcache has been hiding in plain sight; we hope this work motivates others to look inside it.